\documentclass{article}
\usepackage[utf8]{inputenc}
\usepackage{hyperref}
\usepackage{graphicx}
\usepackage{float}
\usepackage[skip=10pt, indent=20pt]{parskip}
\graphicspath{{./figures/}{./lessons/04-digital-io/assets/docs/figures/}}

\title{Lesson 4 - Digital Inputs and Outputs}
\author{Bryan A. \textit{CrazyUncleBurton} Thompson \\ \textit{bryan@batee.com}}
\date{March 13, 2026}

\begin{document}

\maketitle

\section*{Overview}

This lesson will configure the GPIO pins to be digital inputs and outputs.  We 
will also introduce the concept of a port expander to add more inputs and/or 
outputs to the microcontroller.

\section*{Theory}


\section*{Hardware Required}

The following hardware is needed for this lesson:
1 - Tab5 Dev Board
1 - Dual LED Module
1 - Dual Button Module
1 - 4 wire Grove sensor cable
1 - Adafruit 5346 MCP23017 Port Expander w/pins soldered on the top or bottom
1 - Adafruit 4424 QWIIC to Grove cable 
1 - USB cable to connect the Tab5 USB-C port to your PC
1 - Multimeter

\vspace{3mm}


\section*{Objectives}

\begin{itemize}
  \item Learn which pins can be digital inputs and/or outputs
  \item Read Data with a Digital Input
  \item Debounce Switch Contacts
  \item Read Analog Signal with Digital Input 
  \item Write (Output) Data with a Digital Output 
  \item Use a Port Expander for more Digital Input and Output pins
\end{itemize}
\vspace{3mm}



\section*{GPIO Information}
Electrical specs for our GPIO, 
\begin{itemize}
  \item 

\end{itemize}
\vspace{3mm}



\section*{Digital Inputs}
configuring GPIO as digital input, configuring pull-up and pull-down 
resistors, reading a pin vs reading a port vs reading a  reading analog signal as digital input, debouncing inputs

\begin{itemize}
  \item 

\end{itemize}
\vspace{3mm}



\section*{Digital Outputs}

output types (totem pole, open drain, ?!?), configuring GPIO as digital output, 
single bit vs port output vs GPIO register, toggling output,

\begin{itemize}
  \item 

\end{itemize}
\vspace{3mm}



\section*{Port Expanders}
In this section we will add 16 more pins which can be digital inputs or 
outputs, configuring pullups and pull downs, 
\begin{itemize}
  \item 

\end{itemize}
\vspace{3mm}



\section*{Conclusion}
This completes the Digital Input and Output lesson.  In the next lesson, 
we will cover Serial Communication via UART, I2C and SPI!

\end{document}
